\section*{Bab \ref{ch:b3:conformal} --- Pemetaan Konformal dan
Teorema Pemetaan Riemann}

\begin{solution}{exo:b3:conformal:1}
(a) Pemetaan $\varphi$ dan $\psi(w) = \iu\frac{1+w}{1-w}$ berkomposisi menjadi
identitasnya dalam kedua urutannya (lewat perhitungan langsung); sedangkan $\varphi$
memetakan $\mathbb H$ ke dalam $\mathbb D$ dan $\psi$ mengembalikannya
(\cref{ex:b3:conformal:mobius}). Nilainya: $\varphi(\iu) = 0$,
$\varphi(0) = -1$, $\varphi(1) = \frac{1 - \iu}{1 + \iu} =
-\iu$, dan $\varphi(z) \to 1$ saat $z \to \infty$.

(b) Lingkaran dan garis merupakan himpunan penyelesaian
$\alpha\abs z^2 + \bar\beta z + \beta\bar z + \gamma = 0$
(dengan $\alpha, \gamma \in \R$, $\beta \in \C$, dan $\abs\beta^2 >
\alpha\gamma$): sebab $\alpha \neq 0$ memberi lingkaran dan $\alpha = 0$ garis.
Lalu dengan menyubstitusikan $z = 1/w$ dan mengalikannya dengan $\abs w^2$:
$\gamma\abs w^2 + \beta w + \bar\beta\bar w + \alpha = 0$ ---
yakni keluarga yang sama. Sedangkan pemetaan afin jelas mengawetkan keluarganya, dan
setiap pemetaan Möbius merupakan komposisi pemetaan afin dan satu
inversi (sebab $\frac{az+b}{cz+d} = \frac ac + \frac{bc -
ad}{c}\cdot\frac1{cz + d}$ untuk $c \neq 0$).
\end{solution}

\begin{solution}{exo:b3:conformal:2}
(a) Schwarz memberikan $\abs{f(a)} \leq \abs a$ dengan kesamaannya (sebab kedua
ruasnya $= \abs a$): sehingga kasus kesamaannya memaksa $f(z) =
\eu^{\iu\theta}z$, dan $f(a) = a$ memakukan $\eu^{\iu\theta} = 1$.

(b) Misalkan $a \neq b$ titik tetapnya dan $g =
\varphi_a\circ f\circ\varphi_{-a} \in
\operatorname{Aut}(\mathbb D)$ --- dengan memakai
\cref{thm:b3:conformal:autdisc} untuk $\varphi_{\pm a}$. Maka
$g(0) = \varphi_a(f(a)) = 0$ dan $g(c) = c$ untuk $c =
\varphi_a(b) \neq 0$: sehingga menurut (a), $g = \mathrm{id}$, jadi $f =
\varphi_{-a}\circ\varphi_a = \mathrm{id}$.
\end{solution}

\begin{solution}{exo:b3:conformal:3}
Tetapkan $z$ lalu tetapkan $g = \varphi_{f(z)}\circ f\circ\varphi_{-z}$:
yakni pemetaan \omterm{def:b3:holomorphic:holo}{holomorfik} $\mathbb D \to \mathbb D$ dengan $g(0) = 0$, sehingga
$\abs{g'(0)} \leq 1$ (menurut Schwarz). Lalu aturan rantai dengan
$\varphi_a'(\zeta) = \frac{1 - \abs a^2}{(1 - \bar
a\zeta)^2}$:
\[
g'(0) = \varphi_{f(z)}'\bigl(f(z)\bigr)\cdot f'(z)\cdot
\varphi_{-z}'(0)
= \frac{1}{1 - \abs{f(z)}^2}\cdot f'(z)\cdot(1 - \abs z^2),
\]
sehingga diperolehlah ketaksamaan Schwarz--Pick. Sedangkan kesamaan di suatu $z$
membuat $g$ menjadi rotasi, jadi $f =
\varphi_{-f(z)}\circ(\text{rotasi})\circ\varphi_z \in
\operatorname{Aut}(\mathbb D)$ --- lalu kesamaannya berlaku
di mana-mana (hitunglah, atau terapkan ulang dengan peran $f, f^{-1}$
yang ditukar). Jadi pemetaan diri \omterm{def:b3:holomorphic:holo}{holomorfik} cakramnya bersifat
Lipschitz-$1$ terhadap metrik hiperbolik $\frac{2\abs{\dd
z}}{1 - \abs z^2}$; sedangkan automorfismanya adalah isometrinya.
\end{solution}

\begin{solution}{exo:b3:conformal:4}
(a) Pemetaan $z \mapsto \eu^z$: memetakan $\{0 < \operatorname{Im}z <
\pi\}$ secara bijektif ke $\mathbb H$ (sebab $\eu^{x+\iu y} =
\eu^x\eu^{\iu y}$: dengan modulus bebas dan argumen $y \in
\intoo0\pi$), \omterm{def:b3:holomorphic:holo}{holomorfik} dengan turunan yang tak lenyap dan
balikan yang \omterm{def:b3:holomorphic:holo}{holomorfik} (yakni $\log$ pokoknya).
(b) Pemetaan $z \mapsto z^2$ menggandakan argumennya: sehingga kuadran terbuka
$\{0 < \arg z < \frac\pi2\}$ terpetakan secara konformal ke $\mathbb
H$ (dengan balikannya: akar kuadrat pokoknya).
(c) Pemetaan $z \mapsto \frac{1 + z}{1 - z}$ memetakan $\mathbb D$ ke
setengah bidang kanannya dan mengawetkan kesetangkupan atas dan bawahnya: sehingga ia
mengirim setengah cakram atasnya ke kuadran pertama; lalu kuadratkan,
menurut (b), untuk mencapai $\mathbb H$: yakni $z \mapsto \bigl(\frac{1
+ z}{1 - z}\bigr)^2$.
(d) \omterm{thm:b3:conformal:autdisc}{Faktor Blaschke} $\varphi_{1/2}(z) = \frac{z -
\frac12}{1 - \frac z2}$: sebab $\varphi_{1/2}(\tfrac12) = 0$ dan
$\varphi_{1/2}'(\tfrac12) = \frac{1 - \frac14}{(1 -
\frac14)^2} = \frac43 > 0$.
\end{solution}

\begin{solution}{exo:b3:conformal:5}
(a) Sebuah \omterm{def:b3:conformal:conformal}{pemetaan konformal} $\C \to \mathbb D$ (atau $\mathbb H$, setelah
dikomposisikan dengan Cayley) merupakan fungsi utuh yang terbatas:
jadi konstan menurut Liouville --- sehingga tak bijektif. Sedangkan \omterm{def:b3:conformal:conformal}{pemetaan konformal} $\mathbb
D \to \C$ akan berbalikan konformal $\C \to \mathbb D$:
yakni pertentangan yang sama.

(b) Bujur sangkarnya cembung, jadi \omterm{def:b3:conformal:simplyconnected}{terhubung sederhana}, dan sejati:
sehingga secara konformal $\mathbb D$ (\cref{thm:b3:conformal:rmt}). Sedangkan bidang
teriris $\C\setminus\intoc{-\infty}0$ berbentuk bintang terhadap $1$
(sebab ruas dari $1$ menghindari irisannya), jadi \omterm{def:b3:conformal:simplyconnected}{terhubung sederhana},
dan sejati: sehingga secara konformal $\mathbb D$. Untuk cakram tertusuknya: jika
$g \colon \mathbb D\setminus\{0\} \to \mathbb D$ bersifat
konformal, maka $g$ terbatas, sehingga $0$ terhapuskan
(\cref{thm:b3:residues:riemanncw}): jadi $g$ diperluas menjadi $\tilde g
\colon \mathbb D \to \mathbb D$, dan $\tilde g(0)$, karena berada di
peta terbukanya $g(\mathbb D\setminus\{0\}) = \mathbb D$, juga
berupa $g(w)$ untuk suatu $w \neq 0$; lalu dua \omterm{def:b3:topology:topology}{persekitaran} yang saling lepas
atas $0$ dan $w$ berpeta terbuka yang berbagi nilai
$\tilde g(0)$, sehingga berbagi nilai \emph{lain} pula (sebab himpunannya
terbuka): jadi $g$ mengambil suatu nilai dua kali pada
$\mathbb D\setminus\{0\}$ --- yang bertentangan dengan keinjektifannya. Sedangkan untuk
anulus $A = \{1 < \abs z < 2\}$: andaikan $F \colon \mathbb D
\to A$ konformal. Maka $F$ bebas akar pada $\mathbb D$ yang \omterm{def:b3:conformal:simplyconnected}{terhubung
sederhana}, sehingga $F = \eu^L$ untuk $L$ yang \omterm{def:b3:holomorphic:holo}{holomorfik}
(\cref{def:b3:conformal:simplyconnected}). Lalu misalkan $\sigma$
lingkaran $\abs z = \frac32$ di $A$ dan $\gamma =
F^{-1}\circ\sigma$, yakni \omterm{def:b3:topology:pathconnected}{lintasan} tertutup di $\mathbb D$; maka
\[
1 = \operatorname{Ind}_\sigma(0)
= \frac1{2\iu\pi}\int_{F\circ\gamma}\frac{\dd w}{w}
= \frac1{2\iu\pi}\int_\gamma\frac{F'}{F}
= \frac1{2\iu\pi}\int_\gamma L' = 0
\]
(sebab $L'$ berantiturunan): jadi bertentangan. Sehingga baik cakram
tertusuknya maupun anulusnya bukanlah cakram yang menyamar.
\end{solution}

\begin{solution}{exo:b3:conformal:6}
(a) Pemetaan $T(w) = \frac{w - 1}{w + 1}$ memetakan $\{\operatorname{Re}w >
0\}$ secara konformal ke $\mathbb D$ (yakni Cayley yang dirotasikan: sebab $\abs{w -
1} < \abs{w + 1}$ jika dan hanya jika $\operatorname{Re}w > 0$), dengan $T(1) =
0$. Lalu untuk $f \in \mathcal F$, $g = T\circ f\colon \mathbb D \to
\mathbb D$ bersifat \omterm{def:b3:holomorphic:holo}{holomorfik} dengan $g(0) = 0$: sehingga Schwarz memberikan
$\abs{g(z)} \leq \abs z$.

(b) Dengan membalikkan $T$: $f = \frac{1 + g}{1 - g}$, sehingga
\[
\abs{f(z)} \leq \frac{1 + \abs{g(z)}}{1 - \abs{g(z)}}
\leq \frac{1 + \abs z}{1 - \abs z} :
\]
sehingga terbatas lokal (secara seragam pada $\abs z \leq r < 1$). Sedangkan kesamaan
di $z_0 \neq 0$ memaksa $\abs{g(z_0)} = \abs{z_0}$ beserta
kesejajarannya: sehingga $g$ menjadi rotasi, yakni $f(z) = \frac{1 +
\eu^{\iu\theta}z}{1 - \eu^{\iu\theta}z}$ --- yakni
\emph{ekstremal Herglotz}, yaitu \omterm{def:b3:conformal:conformal}{pemetaan konformal} ke setengah bidang
kanannya.
\end{solution}

\begin{solution}{exo:b3:conformal:7}
(i) Keterhubungan sederhananya masuk tepat dua kali, lewat
keberadaan akar kuadrat \omterm{def:b3:holomorphic:holo}{holomorfik} fungsi yang bebas akar
(\cref{def:b3:conformal:simplyconnected}): yakni pada Langkah 0 (untuk akar
$z - b$) dan pada Langkah 2 (untuk akar $\varphi_a\circ f$).
(ii) Syarat $\Omega \neq \C$ menyediakan titik $b$ pada Langkah 0 ---
sebab tanpanya keluarga $\mathcal F$ kosong dari pemetaan terbatas yang
injektif (menurut Liouville). (iii) Sedangkan keterhubungannya dipakai setiap kali
teorema keidentikan atau Hurwitz (\cref{exo:b3:residues:8})
berbicara: sebab limit ekstremalnya ``injektif atau konstan'', dan
kekonstanannya dikecualikan oleh $s > 0$; dan juga pada ``turunan nol
mengakibatkan konstan''. Untuk kegagalannya bagi $\C$: Langkah 0 mustahil, dan
kesimpulannya salah (\cref{exo:b3:conformal:5}(a)).
Sedangkan untuk kegagalannya bagi anulusnya: ia tak \omterm{def:b3:conformal:simplyconnected}{terhubung sederhana} --- sehingga
konstruksi akar kuadratnya patah (misalnya $z$ sendiri, yang bebas akar
pada $A$, tak berakar kuadrat \omterm{def:b3:holomorphic:holo}{holomorfik}: lewat perhitungan indeks
yang sama seperti pada \cref{exo:b3:conformal:5}(b) dengan
$\frac12\int_\sigma\frac{\dd z}z \notin 2\iu\pi\Z$) --- dan
kesimpulannya pun salah.
\end{solution}

\begin{solution}{exo:b3:conformal:8}
Dengan menyubstitusikan uraian Fourier $g$ ke dalam integral
Poissonnya lalu memakai
$\frac1{2\pi}\int P_r(\varphi - t)\eu^{\iu nt}\dd t =
r^{\abs n}\eu^{\iu n\varphi}$ (yang terbaca dari deret $P_r$):
$u(r\eu^{\iu\varphi}) = \sum_nc_n(g)\,r^{\abs
n}\eu^{\iu n\varphi}$, dengan penukarannya dibenarkan kekonvergenan
normalnya (sebab $\abs{c_n} \leq \norm g_\infty$ dan $r < 1$).

(a) Untuk $g = \cos t$: $c_{\pm1} = \frac12$, sehingga $u =
r\cos\varphi = \operatorname{Re}z = x$ --- yang memang harmonik
dengan nilai batas yang tepat.

(b) Karena $\cos^2t = \frac12 + \frac{\cos 2t}2$: maka $u = \frac12 +
\frac{r^2\cos2\varphi}2 = \frac12 +
\frac{\operatorname{Re}(z^2)}2 = \frac12 + \frac{x^2 -
y^2}{2}$.

(c) Untuk $g = \mathbf 1_{(0,\pi)}$ (yakni setengah lingkaran atasnya): $c_0 =
\frac12$ dan $c_n = \frac{1 - (-1)^n}{2\iu\pi n}$ untuk $n \neq
0$, sehingga
\[
u(r\eu^{\iu\varphi}) = \frac12 + \frac2\pi\sum_{k\geq0}
\frac{r^{2k+1}\sin\bigl((2k+1)\varphi\bigr)}{2k + 1},
\qquad u(0) = \frac12 :
\]
sehingga pusatnya melihat tepat rata-rata data batasnya ---
yakni sifat nilai rata-rata yang menjelma.
\end{solution}

\begin{solution}{exo:b3:conformal:9}
Dari $(1-r)^2 \leq 1 - 2r\cos\theta + r^2 \leq (1+r)^2$:
\[
\frac{1-r}{1+r} = \frac{1 - r^2}{(1+r)^2} \leq P_r(\theta)
\leq \frac{1 - r^2}{(1 - r)^2} = \frac{1+r}{1-r} .
\]
Lalu untuk $u \geq 0$ yang harmonik pada $\mathbb D$ dan $s < 1$: $u_s(z) =
u(sz)$ bersifat harmonik pada sebuah \omterm{def:b3:topology:topology}{persekitaran} $\bar{\mathbb D}$,
sehingga sama dengan integral Poissonnya
(\cref{thm:b3:conformal:poisson}, lewat ketunggalannya, yang diterapkan pada
nilai batasnya sendiri); lalu dengan menjepit kernelnya dan memakai
nilai rata-ratanya $\frac1{2\pi}\int u_s(\eu^{\iu t})\dd t = u(0)$:
\[
\frac{1-r}{1+r}\,u(0) \leq u(s\,r\eu^{\iu\varphi}) \leq
\frac{1+r}{1-r}\,u(0).
\]
Lalu biarkan $s \to 1^-$ pada $r\eu^{\iu\varphi}$ yang tetap (lewat \omterm{def:b3:topology:continuity}{kekontinuan}
$u$): diperolehlah ketaksamaan Harnack. Sedangkan jika $u$ harmonik pada $\C$
dengan $u \geq m$: terapkan Harnack pada $u - m$ di cakram $D(0, R)$,
yakni pada $z \mapsto u(Rz) - m$: maka untuk $z$ yang tetap dan $r =
\abs z/R \to 0$, kedua batasnya menuju $u(0) - m$: sehingga $u(z) =
u(0)$ --- jadi konstan (yakni Liouville dua sisi dari sebuah batas
sepihak).
\end{solution}

\begin{solution}{exo:b3:conformal:10}
Secara lokal, $u = \operatorname{Re}F$ dengan $F$ yang \omterm{def:b3:holomorphic:holo}{holomorfik}
(\cref{prop:b3:conformal:harmonicholo}), sehingga $u\circ f =
\operatorname{Re}(F\circ f)$ bersifat harmonik di mana pun ia terdefinisi:
jadi keharmonikannya invarian secara konformal. Pada $\mathbb H$: logaritma
pokoknya memberikan $\log z = \ln\abs z + \iu\arg z$ yang
\omterm{def:b3:holomorphic:holo}{holomorfik} pada $\mathbb H$, sehingga $u = \frac1\pi\arg z =
\operatorname{Im}\bigl(\frac1\pi\log z\bigr)$ bersifat harmonik,
dengan limit batasnya: sebab untuk $x > 0$, $\arg \to 0$ dan $u \to 0$;
sedangkan untuk $x < 0$, $\arg \to \pi$ dan $u \to 1$: yakni data $\mathbf
1_{\intoo{-\infty}0}$ di setiap $x \neq 0$. Lalu dengan mengangkutnya lewat
\omterm{ex:b3:conformal:mobius}{pemetaan Cayley} (yang mengirim $\mathbb D \to \mathbb H$ setelah
inversinya dan mencocokkan setengah lingkaran atasnya dengan sumbu
negatifnya, sampai rotasi yang ditetapkan lewat pengejaran tiga titik
batasnya), $u\circ(\text{Cayley})$ menyelesaikan soal cakram
\cref{exo:b3:conformal:8}(c); lalu menilaikannya di pusatnya
memulihkan $u = \frac12$ di sana, dan bentuk tertutupnya
$\frac1\pi\arg$ dapat diperiksa terhadap deretnya dengan menjumlahkan
identitas bertipe $\sum\frac{r^{2k+1}\sin((2k+1)\varphi)}{2k+1} =
\frac12\arctan\frac{2r\sin\varphi}{1 - r^2}$
--- yakni jalur dasar menuju jawaban yang sama.
\end{solution}

\begin{solution}{exo:b3:conformal:11}
(a) Misalkan $a \neq b$ tetap. Lalu dengan mengonjugatkannya lewat $\varphi_a(z) =
\frac{z - a}{1 - \bar az}$ (yakni automorfisma yang menukar $a$
dan $0$), $g = \varphi_a\circ f\circ\varphi_a^{-1}$ menetapkan
$0$ dan titik $c = \varphi_a(b) \neq 0$. Lalu Schwarz:
$\abs{g(z)} \leq \abs z$, dan di $z = c$ kesamaannya berlaku
(sebab $g(c) = c$): sehingga kasus kesamaannya memaksa $g(z) = \lambda z$
dengan $\abs\lambda = 1$, dan $\lambda c = c$ memberikan $\lambda =
1$: jadi $g = \mathrm{id}$, sehingga $f = \mathrm{id}$.

(b) Fungsi $h(z) = f(z)/z$ (dengan \omterm{thm:b3:residues:riemanncw}{kesingularan terhapuskan} di $0$) bersifat
\omterm{def:b3:holomorphic:holo}{holomorfik} pada $\mathbb D$ dengan $\abs h \leq 1$ (menurut Schwarz);
dan $\abs h < 1$ di mana-mana, sebab jika tidak asas maksimumnya
(lewat maksimum \omterm{def:b3:topology:interior}{interior} $\abs h$) akan membuat $h$ menjadi konstanta
unimodular, yakni $f$ sebuah rotasi --- yang terkecualikan. Lalu pada \omterm{def:b3:topology:compact}{kompak}
$\bar D(0,r)$, $c_r = \max\abs h < 1$, sehingga $\abs{f(z)} \leq
c_r\abs z$ di sana; lebih lanjut $f$ memetakan $\bar D(0,r)$ ke
dirinya (sebab $c_r\abs z \leq r$), sehingga batasnya teriterasi:
$\abs{f^{\circ n}(z)} \leq c_r^n\,r \to 0$ secara seragam pada
$\bar D(0, r)$.

(c) Titik tetap $\frac{z^2 + z}2$: sebab $z^2 + z = 2z$ jika dan hanya jika
$z(z - 1) = 0$; sehingga hanya $z = 0$ yang terletak di $\mathbb D$ (sebab $z = 1$
pada batasnya). Ia bukan rotasi (sebab $f'(0) = \frac12$), sehingga
\omterm{def:b3:groups:action}{orbitnya} menuju $0$; dan secara kuantitatif $f(z) = \frac z2(1 + z)$
memberikan $\abs{f(z)} \leq \frac{3}{4}\abs z$ pada $\abs z \leq
\frac12$, lalu begitu \omterm{def:b3:groups:action}{orbitnya} kecil, $\abs{f(z)} \approx
\frac{\abs z}2$: jadi asimtotik geometri bernisbah
$f'(0) = \frac12$. Dari $z_0 = \frac12$: $z_1 = \frac38$,
$z_2 \approx 0.258$, dan $z_3 \approx 0.162$ --- yakni separuh setiap
langkahnya, seperti yang diramalkan pengalinya.
\end{solution}

\begin{solution}{exo:b3:conformal:12}
(a) Berlaku $\Delta u = 6x - 6x + 0 = 0$. Lalu Cauchy--Riemann menuntut
$v_y = u_x = 3x^2 - 3y^2$ dan $v_x = -u_y = 6xy - 2$.
Dengan mengintegralkan yang pertama terhadap $y$: $v = 3x^2y - y^3 + c(x)$;
lalu memasukkannya ke yang kedua: $6xy + c'(x) = 6xy - 2$, sehingga $c(x)
= -2x + C$. Jadi $v = 3x^2y - y^3 - 2x + C$ dan
\[
f = u + \iu v = (x^3 - 3xy^2) + \iu(3x^2y - y^3)
+ 2y - 2\iu x + \iu C = z^3 - 2\iu z + \iu C .
\]

(b) Bentuk $\omega = -u_y\,\dd x + u_x\,\dd y$ bersifat tertutup
justru karena $\Delta u = 0$ (sebab $\partial_y(-u_y) =
-u_{yy} = u_{xx} = \partial_x(u_x)$). Lalu pada \omterm{def:b3:topology:topology}{himpunan terbuka} berbentuk
bintang lema Poincaré (\cref{thm:b3:forms:poincare}; atau
konstruksi antiturunan
\cref{thm:b3:holomorphic:cauchy} yang diterapkan pada
$u_x - \iu u_y$ yang \omterm{def:b3:holomorphic:holo}{holomorfik}) menyediakan $v$ dengan $\dd v = \omega$, yakni
sistem Cauchy--Riemannya: sehingga $u + \iu v$ \omterm{def:b3:holomorphic:holo}{holomorfik}. Sedangkan dua
sekawannya berselisih sebuah fungsi yang gradiennya lenyap: yakni sebuah
konstanta (lewat keterhubungannya).

(c) Untuk $u = \ln\abs z$ pada $\C^*$: $\omega = -u_y\dd x +
u_x\dd y = \frac{-y\,\dd x + x\,\dd y}{x^2 + y^2} =
\omega_\theta$, yakni bentuk sudutnya
(\cref{ex:b3:forms:angular}), yang integralnya sepanjang lingkaran
satuannya adalah $2\pi \neq 0$: jadi tak eksak, sehingga tak ada sekawan
global --- sebab sebuah sekawan akan berupa penentuan
argumen yang \omterm{def:b3:topology:continuity}{kontinu}, dan bilangan lilitannya tepat merupakan
rintangannya. Sedangkan secara lokal (pada sembarang subdaerah berbentuk bintang), $v =
\arg z$ berhasil dan $u + \iu v = \log z$: jadi kegagalannya
global, bukan lokal.
\end{solution}

\begin{solution}{pb:b3:conformal:1}
\textbf{1.} Fungsi $g$ bersifat injektif dan \omterm{def:b3:holomorphic:holo}{holomorfik}; sedangkan pada $C_\rho$
(dengan $\rho > 1$) ia mulus, dan luas terlingkupinya adalah
$\frac1{2\iu}\oint_{g(C_\rho)}\bar\zeta\,\dd\zeta$ (yakni rumus luas
Green--Riemann Tahun ke-2, yang diterapkan dengan arah positif).
Lalu dengan menyubstitusikan $\zeta = g(w)$ dan $w =
\rho\eu^{\iu\theta}$:
\[
A_\rho = \frac1{2\iu}\int_{C_\rho}\overline{g(w)}\,g'(w)\dd
w = \frac\rho2\int_0^{2\pi}\overline{g(\rho\eu^{\iu\theta})}
\,g'(\rho\eu^{\iu\theta})\,\eu^{\iu\theta}\,\dd\theta .
\]
Sisipkan $\bar g = \rho\eu^{-\iu\theta} + \bar b_0 +
\sum_m\bar b_m\rho^{-m}\eu^{\iu m\theta}$ dan $g' = 1 -
\sum_nnb_n\rho^{-n-1}\eu^{-\iu(n+1)\theta}$: maka setelah
dikalikan $\eu^{\iu\theta}$, hanya hasil kali berfrekuensi nol yang
bertahan pada pengintegralan-$\theta$-nya (sebab kekonvergenan
normalnya membenarkan kerja suku demi sukunya): sehingga pasangan
$(\rho\eu^{-\iu\theta})\cdot1$ menyumbang $2\pi\rho$, dan
setiap pasangan $\bar b_n\rho^{-n}\eu^{\iu
n\theta}\cdot(-nb_n\rho^{-n-1}\eu^{-\iu(n+1)\theta})$
menyumbang $-2\pi n\abs{b_n}^2\rho^{-2n-1}$:
\[
A_\rho = \pi\rho^2 - \pi\sum_{n\geq1}n\abs{b_n}^2\rho^{-2n}.
\]

\textbf{2.} Luas bersifat tak negatif: sehingga $\sum_{n\leq
N}n\abs{b_n}^2\rho^{-2n} \leq \rho^2$ untuk setiap $N$; lalu biarkan
$\rho \to 1^+$ kemudian $N \to \infty$: diperoleh
$\sum_{n\geq1}n\abs{b_n}^2 \leq 1$. Sedangkan kesamaan pada $\abs{b_1}
\leq 1$ memaksa semua $b_n = 0$ lainnya: jadi $g(w) = w + b_0 +
\eu^{\iu\alpha}/w$, yang memetakan ke komplemen sebuah
ruas berpanjang $4$ (yakni pemetaan bertipe Joukowski): itulah ekstremalnya.

\textbf{3.} Fungsi $f(z)/z = 1 + a_2z + \cdots$ bersifat \omterm{def:b3:holomorphic:holo}{holomorfik} dan
bebas akar pada $\mathbb D$ (sebab $f$ lenyap hanya di $0$, secara sederhana:
menurut keinjektifannya), sehingga demikian pula $f(z^2)/z^2$, yang berakar
kuadrat \omterm{def:b3:holomorphic:holo}{holomorfik} $\varphi$ dengan $\varphi(0) = 1$
(\cref{def:b3:conformal:simplyconnected}; sebab $\mathbb D$
cembung). Maka $h(z) = z\varphi(z^2)$ memenuhi $h^2 =
f(z^2)$, $h(z) = z\bigl(1 + \frac{a_2}2z^2 + \cdots\bigr)$
(lewat deret binomial akarnya), dan $h(z) = z\varphi(z^2)$ bersifat
gasal menurut konstruksinya (sebab $\varphi(z^2)$ genap). Untuk keinjektifannya: $h(z) = h(z')$ memberikan $f(z^2) =
f(z'^2)$, sehingga $z^2 = z'^2$, yakni $z' = \pm z$; lalu jika $z' = -z$
maka kegasalannya memberikan $h(z) = -h(z)$, sehingga $h(z) = 0$, yang memaksa $z
= 0$ (sebab $\varphi$ bebas akar): jadi $z = z' = 0$.

\textbf{4.} Untuk $g(w) = 1/h(1/w)$: bagi $\abs w > 1$, $1/w \in
\mathbb D\setminus\{0\}$ dan $h \neq 0$ di sana: jadi terdefinisi baik,
injektif (sebab komposisi injeksi), dengan uraian
\[
g(w) = \frac{1}{\frac1w\bigl(1 + \frac{a_2}{2w^2} +
\cdots\bigr)} = w\Bigl(1 - \frac{a_2}{2w^2} + \cdots\Bigr)
= w - \frac{a_2}{2}\,w^{-1} + \cdots :
\]
yakni berbentuk Bagian I dengan $b_1 = -\frac{a_2}2$. Lalu teorema luasnya
memberikan $\abs{\frac{a_2}2} \leq 1$: sehingga $\abs{a_2} \leq 2$.

\textbf{5.} Berlaku $k(z) = \frac z{(1-z)^2} = z\sum_{m\geq0}(m +
1)z^m = \sum_{n\geq1}nz^n$: sehingga koefisiennya $a_n = n$, jadi $a_2 =
2$. Untuk keunivalenan dan petanya: $k = \frac14\bigl[w^2 - 1\bigr]$
dengan $w = \frac{1 + z}{1 - z}$, yakni \omterm{def:b3:conformal:conformal}{pemetaan konformal} $\mathbb D$
ke setengah bidang kanannya; lalu $w^2$ memetakan setengah bidang itu
secara konformal ke $\C\setminus\intoc{-\infty}0$; kemudian
$\frac{\cdot - 1}4$ memberikan
$\C\setminus\intoc{-\infty}{-\frac14}$: jadi injektif pada setiap
tahapnya, dengan peta seperti yang diklaim.

\textbf{6.} Untuk $F = \frac{cf}{c - f}$: karena $c \notin f(\mathbb
D)$, penyebutnya tak pernah lenyap: sehingga $F$ \omterm{def:b3:holomorphic:holo}{holomorfik}, dan
injektif (sebab $w \mapsto \frac{cw}{c - w}$ bersifat Möbius dan injektif
di luar $w = c$). Untuk uraiannya: dengan $f = z + a_2z^2 + \cdots$,
\[
F = f\cdot\frac1{1 - f/c} = \bigl(z + a_2z^2\bigr)\Bigl(1 +
\frac zc\Bigr) + O(z^3)
= z + \Bigl(a_2 + \frac1c\Bigr)z^2 + O(z^3):
\]
sehingga $F \in \mathcal S$ dengan $A_2 = a_2 + \frac1c$.

\textbf{7.} Bieberbach dua kali: $\abs{a_2} \leq 2$ dan
$\abs{a_2 + \frac1c} \leq 2$, sehingga $\abs{\frac1c} \leq 4$:
jadi $\abs c \geq \frac14$. Sehingga setiap nilai yang dilewatkan terletak di luar
$D(0,\frac14)$, yakni $f(\mathbb D) \supseteq D(0, \frac14)$.
Dan tajam: sebab fungsi Koebenya melewatkan $-\frac14$ (pertanyaan 5).

\textbf{8.} Misalkan $d = d(z_0, \partial\Omega)$ dan $\psi =
\varphi^{-1} \colon \mathbb D \to \Omega$ dengan $\psi(0) = z_0$ dan
$\psi'(0) = 1/\varphi'(z_0) > 0$. Maka penormalannya $\tilde
f(z) = \frac{\psi(z) - z_0}{\psi'(0)}$ terletak di $\mathcal S$,
sehingga petanya memuat $D(0, \frac14)$; lalu dengan menskalakannya kembali,
$\Omega = \psi(\mathbb D) \supseteq D\bigl(z_0,
\tfrac{\psi'(0)}4\bigr)$, sehingga $d \geq \frac{\psi'(0)}4 =
\frac1{4\varphi'(z_0)}$: yakni ketaksamaan kirinya. Sedangkan untuk yang kanan: $\varphi\circ(z_0 + d\,\cdot)$
memetakan $\mathbb D$ ke dalam $\mathbb D$ dengan $0 \mapsto 0$; sehingga Schwarz
membatasi turunannya di $0$: $d\,\varphi'(z_0) \leq 1$, yakni
$\frac1{\varphi'(z_0)} \geq d$ --- sehingga bersama-sama
\[
d \leq \frac1{\varphi'(z_0)} \leq 4d .
\]
(Sedangkan ketaksamaan kiri seperti yang ditampilkan pada pernyataannya adalah
rantai yang sama, yang disusun ulang.)

\textbf{9.} Untuk fungsi Koebe $a_n = n$: sehingga kesamaannya berlaku
di sepanjang dugaannya; dan rotasinya
$\eu^{-\iu\theta}k(\eu^{\iu\theta}z) = \sum
n\eu^{\iu(n-1)\theta}z^n$ mempunyai $\abs{a_n} = n$ pula. Lalu dengan mendorong
pertanyaan 4: $h = z + \frac{a_2}2z^3 + c_5z^5 + \cdots$ dengan
$h^2 = f(z^2)$; sehingga dengan membandingkan koefisien $z^6$-nya: $2c_5 +
\frac{a_2^2}4 = a_3$, jadi $c_5 = \frac{a_3}2 -
\frac{a_2^2}8$. Maka
\[
g(w) = \frac1{h(1/w)} = w\Bigl(1 - \frac{a_2}{2w^2} +
\Bigl(\frac{a_2^2}4 - c_5\Bigr)\frac1{w^4} + \cdots\Bigr)
= w - \frac{a_2}2w^{-1} + \Bigl(\frac{3a_2^2}8 -
\frac{a_3}2\Bigr)w^{-3} + \cdots,
\]
dan teorema luasnya (yakni $\sum n\abs{b_n}^2 \leq 1$) melahirkan
\[
\Bigl|\frac{a_2}2\Bigr|^2 + 3\,\Bigl|\frac{3a_2^2}8 -
\frac{a_3}2\Bigr|^2 \leq 1 .
\]
Untuk pemeriksaan Koebenya (dengan $a_2 = 2$ dan $a_3 = 3$): $\abs{\frac{a_2}2}^2 = 1$
dan $\frac{3\cdot4}8 - \frac32 = 0$: sehingga totalnya tepat $1$ ---
jadi ekstremal, seperti yang seharusnya.

\textbf{10.} Pemetaan $\varphi$ merupakan automorfisma cakram dengan
$\varphi(0) = z_0$, sehingga $f\circ\varphi$ bersifat univalen pada
$\mathbb D$ (sebab komposisi injeksi), dan $f'(z_0) \neq
0$ (\cref{def:b3:conformal:conformal}): jadi $F$ terdefinisi baik
dan univalen, dengan $F(0) = 0$. Lalu lewat aturan hasil bagi:
\[
\varphi'(z) = \frac{(1 + \bar z_0z) - (z + z_0)\bar z_0}
{(1 + \bar z_0z)^2}
= \frac{1 - \abs{z_0}^2}{(1 + \bar z_0z)^2},
\qquad \varphi'(0) = 1 - \abs{z_0}^2 ,
\]
sehingga $(f\circ\varphi)'(0) = f'(z_0)(1 - \abs{z_0}^2)$ dan
$F'(0) = 1$: jadi $F \in \mathcal S$.

\textbf{11.} Dengan $G = f\circ\varphi$: $G'' =
f''(\varphi)\,\varphi'^2 + f'(\varphi)\,\varphi''$, dan
$\varphi''(z) = -2\bar z_0(1 - \abs{z_0}^2)(1 + \bar
z_0z)^{-3}$ memberikan $\varphi''(0) = -2\bar z_0(1 -
\abs{z_0}^2)$. Karena itu
\[
A_2 = \frac{G''(0)}{2f'(z_0)(1 - \abs{z_0}^2)}
= \frac12\Bigl[\bigl(1 - \abs{z_0}^2\bigr)
\frac{f''(z_0)}{f'(z_0)} - 2\bar z_0\Bigr] .
\]
Lalu Bieberbach untuk $F$ (pertanyaan 4): $\abs{A_2} \leq 2$, yakni
$\bigl|(1 - \abs{z_0}^2)\frac{f''(z_0)}{f'(z_0)} - 2\bar
z_0\bigr| \leq 4$. Kalikanlah dengan
$z_0/(1 - \abs{z_0}^2)$, yang bermodulus
$r/(1 - r^2)$ untuk $z_0 = r\eu^{\iu\theta}$, lalu pakailah $z_0\bar
z_0 = r^2$:
\[
\Bigl|\,z_0\frac{f''(z_0)}{f'(z_0)} -
\frac{2r^2}{1 - r^2}\Bigr| \leq \frac{4r}{1 - r^2} .
\]

\textbf{12.} Sebuah bilangan kompleks yang berjarak paling banyak $\rho$ dari
titik real $c$ berbagian real di $\intcc{c - \rho}{c +
\rho}$: lalu terapkan ini pada $c = \frac{2r^2}{1-r^2}$ dan $\rho =
\frac{4r}{1-r^2}$.

\textbf{13.} Untuk fungsi $\mathcal C^1$ tak lenyap $g$:
$\log\abs g = \frac12\log(g\bar g)$, sehingga $\frac{\dd}{\dd
t}\log\abs g = \frac{g'\bar g + g\bar g'}{2\abs g^2} =
\operatorname{Re}\frac{g'}g$. Lalu dengan $g(t) =
f'(t\eu^{\iu\theta})$ (yang bebas akar menurut keunivalenannya), $g'(t) =
\eu^{\iu\theta}f''(t\eu^{\iu\theta})$, sehingga di $z =
t\eu^{\iu\theta}$:
\[
\frac{\dd}{\dd t}\log\bigl|f'(t\eu^{\iu\theta})\bigr|
= \operatorname{Re}\Bigl(\eu^{\iu\theta}
\frac{f''}{f'}\Bigr)
= \frac1t\operatorname{Re}\Bigl(z\frac{f''}{f'}\Bigr)
\in \Bigl[\frac{2t - 4}{1 - t^2},\
\frac{2t + 4}{1 - t^2}\Bigr]
\]
menurut pertanyaan 12 pada jari-jari $t$. Lalu karena $\frac{\dd}{\dd
t}\log\frac{1+t}{(1-t)^3} = \frac1{1+t} + \frac3{1-t} =
\frac{2t+4}{1-t^2}$ dan $\frac{\dd}{\dd
t}\log\frac{1-t}{(1+t)^3} = \frac{-1}{1-t} - \frac3{1+t} =
\frac{2t-4}{1-t^2}$, mengintegralkannya dari $0$ ke $r$ (sebab ketiga
fungsinya lenyap di $t = 0$ dan $f'(0) = 1$) memberikan
\[
\log\frac{1-r}{(1+r)^3} \leq \log\abs{f'(z)} \leq
\log\frac{1+r}{(1-r)^3} :
\]
yakni teorema kesesatannya, setelah dieksponensialkan.

\textbf{14.} Untuk yang atas: sepanjang ruas $[0, z]$,
\[
\abs{f(z)} = \Bigl|\int_0^rf'(t\eu^{\iu\theta})
\eu^{\iu\theta}\dd t\Bigr|
\leq \int_0^r\frac{1+t}{(1-t)^3}\dd t
= \frac{r}{(1-r)^2}
\]
(sebab $\frac{\dd}{\dd t}\frac t{(1-t)^2} = \frac{1+t}{(1-t)^3}$).
Untuk yang bawah: perhatikan $\frac r{(1+r)^2} < \frac14$ bagi $r < 1$ (sebab ia
menyatakan $(1-r)^2 > 0$), sehingga jika $\abs{f(z)} \geq \frac14$ maka tak ada
yang perlu dibuktikan. Selain itu ruas $[0, f(z)]$ terletak
di $D(0, \frac14) \subseteq f(\mathbb D)$ (pertanyaan 7), dan
$\gamma = f^{-1}\circ[0, f(z)]$ merupakan \omterm{def:b3:topology:pathconnected}{lintasan} $\mathcal C^1$
dari $0$ ke $z$ di $\mathbb D$ (sebab $f^{-1}$ \omterm{def:b3:holomorphic:holo}{holomorfik},
\cref{cor:b3:residues:openmapping}). Lalu dengan menyubstitusikan $w =
f(\zeta)$,
\[
\abs{f(z)} = \int_{[0,f(z)]}\abs{\dd w}
= \int_\gamma\abs{f'(\zeta)}\,\abs{\dd\zeta}
\geq \int_0^1\frac{1 - \rho(s)}{(1 + \rho(s))^3}
\,\abs{\gamma'(s)}\,\dd s
\]
dengan $\rho = \abs\gamma$, memakai batas kesesatan bawahnya
pada jari-jari $\rho(s)$. Lalu karena $\abs{\gamma'} \geq \rho'$
(di mana pun ia terdefinisi; sebab $\rho$ bersifat Lipschitz) dan
faktor integrannya positif,
\[
\abs{f(z)} \geq \int_0^1
\frac{1 - \rho(s)}{(1 + \rho(s))^3}\,\rho'(s)\,\dd s
= \int_0^{r}\frac{1 - u}{(1 + u)^3}\,\dd u
= \frac{r}{(1 + r)^2}
\]
(sebab $\rho(0) = 0$ dan $\rho(1) = r$; dan substitusinya hanya memakai sebuah
antiturunan, $\frac{\dd}{\dd u}\frac u{(1+u)^2} =
\frac{1-u}{(1+u)^3}$, bukan kemonotonannya).

\textbf{15.} Berlaku $k'(z) = \frac{\dd}{\dd z}\,z(1 - z)^{-2} =
\frac{1 + z}{(1 - z)^3}$. Di $z = r$: $k'(r) =
\frac{1+r}{(1-r)^3}$ dan $k(r) = \frac r{(1-r)^2}$ --- sehingga kedua
batas atasnya tercapai. Sedangkan di $z = -r$: $k'(-r) =
\frac{1-r}{(1+r)^3}$ dan $\abs{k(-r)} = \frac r{(1+r)^2}$
--- sehingga kedua batas bawahnya tercapai. Jadi fungsi Koebe merentangkan
sumbu positifnya secara maksimal menuju batas yang jauh dan
mengerutkan sinar antipodalnya secara maksimal ke arah ujung irisannya.

\textbf{16.} Bahwa $\abs{a_2} = 2$ berarti $\abs{b_1} = 1$ bagi $g(w)
= 1/h(1/w) = w - \frac{a_2}2w^{-1} + \cdots$ (pertanyaan 4);
lalu teorema luasnya $\sum n\abs{b_n}^2 \leq 1$ membunuh setiap
koefisien lainnya: sehingga $g(w) = w + b_0 + \eu^{\iu\alpha}/w$ dengan
$\eu^{\iu\alpha} = -\frac{a_2}2$. Lalu karena $h$ gasal, $g(-w) =
1/h(-1/w) = -1/h(1/w) = -g(w)$: sehingga $g$ gasal, jadi $b_0 = 0$.
Lalu dengan membalikkannya, $h(z) = 1/g(1/z) = \frac{z}{1 +
\eu^{\iu\alpha}z^2}$, dan $f(z^2) = h(z)^2$ memberikan
\[
f(w) = \frac{w}{(1 + \eu^{\iu\alpha}w)^2}
= \frac{w}{(1 - \eu^{\iu\theta}w)^2}
= \eu^{-\iu\theta}k\bigl(\eu^{\iu\theta}w\bigr),
\qquad \eu^{\iu\theta} = -\eu^{\iu\alpha} .
\]
Sebaliknya setiap rotasinya mempunyai $a_2 = 2\eu^{\iu\theta}$ yang
bermodulus $2$: sehingga ekstremal Bieberbach tepat berupa
fungsi Koebe yang dirotasikan.

\textbf{17.} Jika $c$ dilewatkan dengan $\abs c = \frac14$:
maka pertanyaan 6 memberikan $F \in \mathcal S$ dengan $A_2 = a_2 +
\frac1c$, sehingga $\frac1c = A_2 - a_2$ dan
\[
4 = \Bigl|\frac1c\Bigr| = \abs{A_2 - a_2}
\leq \abs{A_2} + \abs{a_2} \leq 2 + 2 = 4 :
\]
sehingga kesamaannya di sepanjangnya, khususnya $\abs{a_2} = 2$. Lalu menurut
pertanyaan 16, $f = \eu^{-\iu\theta}k(\eu^{\iu\theta}\cdot)$,
yang himpunan terlewatnya adalah $\eu^{-\iu\theta}
\intoc{-\infty}{-\frac14}$: sehingga satu-satunya nilai terlewat yang
bermodulus $\frac14$ adalah $c = -\eu^{-\iu\theta}/4$. (Periksalah:
$a_2 = 2\eu^{\iu\theta}$ dan $\frac1c = -4\eu^{\iu\theta}$,
sehingga $A_2 = -2\eu^{\iu\theta} = -a_2$: jadi ketaksamaan
segitiganya jenuh lewat pelawanan arah, seperti yang seharusnya.)

\textbf{18.} \omterm{thm:b3:holomorphic:analytic}{Taksiran Cauchy} pada lingkaran $\abs z =
r$ (\cref{thm:b3:holomorphic:analytic}), yang digabungkan dengan
teorema pertumbuhannya, memberikan
\[
\abs{a_n} \leq r^{-n}\sup_{\abs z = r}\abs f
\leq r^{1-n}\,(1 - r)^{-2} .
\] Lalu pilihlah $r = 1 - \frac1n$ (dengan $n \geq
2$): maka $r^{1-n} = \bigl(1 + \frac1{n-1}\bigr)^{n-1} < \eu$
(sebab barisannya naik dengan limit $\eu$) dan $(1 - r)^{-2} =
n^2$: sehingga $\abs{a_n} < \eu\,n^2$.

\textbf{19.} Himpunan $U = f(D(0, r))$ bersifat terbuka
(\cref{cor:b3:residues:openmapping}) dan memuat $0$.
Untuk batasnya: jika $p \in \partial U$, tulislah $p = \lim f(z_k)$
dengan $z_k \in D(0, r)$; maka sebuah subbarisan memberikan $z_k \to z_\infty
\in \bar D(0, r)$, dan $f(z_\infty) = p \notin U$ memaksa
$z_\infty \in \partial D(0, r)$: sehingga $\partial U \subseteq
f(\partial D(0, r))$, jadi setiap titik batas $U$ bermodulus
$\geq \frac r{(1+r)^2}$ (pertanyaan 14). Kini misalkan
$\abs w < \frac r{(1+r)^2}$ dan andaikan $w \notin U$. Maka
ruas $[0, w]$ \omterm{def:b3:topology:connected}{terhubung}, memotong $U$ (di $0$) dan
komplemennya (di $w$), sehingga ia memotong $\partial U$; padahal semua
titiknya bermodulus $\leq \abs w < \frac r{(1+r)^2}$:
jadi bertentangan. Karena itu $D\bigl(0, \frac r{(1+r)^2}\bigr)
\subseteq U$. Untuk ketajamannya: $k(-r) = -\frac r{(1+r)^2}$ merupakan
peta sebuah titik batas $D(0,r)$, dan $k$ injektif,
sehingga $k(-r) \notin k(D(0, r))$: jadi jari-jarinya tak dapat
diperbesar. Lalu saat $r \to 1^-$, $\frac r{(1+r)^2} \to \frac14$:
yakni teorema seperempatnya bagi cakram penuhnya.

\textbf{20.} Pertama, $\partial f(\mathbb D) \neq \emptyset$:
sebab jika tidak, $f(\mathbb D)$, yang terbuka, tertutup, dan tak kosong, akan
seluruh $\C$, sehingga $f^{-1} \colon \C \to \mathbb D$ akan berupa
fungsi utuh yang terbatas dan tak konstan, yang bertentangan dengan
\cref{cor:b3:holomorphic:liouville}. Untuk ketaksamaan kirinya:
transformasi Koebe $F$ pada pertanyaan 10 terletak di $\mathcal S$ dan
$F(\mathbb D) = \frac{f(\mathbb D) -
f(z_0)}{f'(z_0)(1-\abs{z_0}^2)}$; sehingga menurut pertanyaan 7 ia memuat
$D(0, \frac14)$, jadi
\[
f(\mathbb D) \supseteq f(z_0) + D\Bigl(0,\
\tfrac14\abs{f'(z_0)}\bigl(1 - \abs{z_0}^2\bigr)\Bigr) ,
\]
dan setiap titik $\partial f(\mathbb D)$ --- yang saling lepas dengan
\omterm{def:b3:topology:topology}{himpunan terbuka} $f(\mathbb D)$ --- berjarak $\geq
\frac14(1-\abs{z_0}^2)\abs{f'(z_0)}$ dari $f(z_0)$. Untuk ketaksamaan
kanannya: misalkan $d = d(f(z_0), \partial f(\mathbb D)) <
\infty$. Maka cakram $D(f(z_0), d)$ terletak di $f(\mathbb D)$: sebab
sebuah ruas dari $f(z_0)$ ke sembarang titiknya tetap berjarak
$< d$ dari $f(z_0)$, sehingga tak pernah memotong $\partial f(\mathbb D)$,
dan argumen keterhubungan pertanyaan 19 menjaganya di
$f(\mathbb D)$. Lalu $\hat g(w) = f^{-1}(f(z_0) + dw)$ memetakan
$\mathbb D$ ke dalam $\mathbb D$ dengan $\hat g(0) = z_0$, dan
$\chi = \psi^{-1}\circ\hat g$, dengan $\psi(z) = \frac{z +
z_0}{1 + \bar z_0z}$, menetapkan $0$: sehingga Schwarz
(\cref{thm:b3:conformal:schwarz}) memberikan $\abs{\chi'(0)}
\leq 1$. Lalu karena $\chi'(0) = \frac{\hat g'(0)}{\psi'(0)} =
\frac{d}{f'(z_0)\,(1 - \abs{z_0}^2)}$, ini berarti $d \leq (1 -
\abs{z_0}^2)\abs{f'(z_0)}$.

\textbf{21.} Berlaku $k(\mathbb D) = \C\setminus
\intoc{-\infty}{-\frac14}$, sehingga $\partial k(\mathbb D) =
\intoc{-\infty}{-\frac14}$, dan untuk titik real positif
$k(r) = \frac r{(1-r)^2}$ titik batas terdekatnya adalah
$-\frac14$:
\[
d = k(r) + \frac14 = \frac{4r + (1-r)^2}{4(1-r)^2}
= \frac{(1+r)^2}{4(1-r)^2} .
\]
Sedangkan anggota kiri pertanyaan 20: $\frac14(1 - r^2)k'(r) =
\frac14\,\frac{(1-r^2)(1+r)}{(1-r)^3} =
\frac{(1+r)^2}{4(1-r)^2} = d$: yakni kesamaan. (Sedangkan anggota kanannya
bernilai $4d$: sehingga faktor $4$ penuhnya memisahkan kedua sisinya,
dan fungsi Koebe duduk tepat di dasarnya.)

\textbf{22.} Asas tunggalnya adalah \omterm{thm:b3:hilbert:parseval}{Parseval}: sebab keunivalenannya
melarang pertindihan, sehingga luas komplemen peta
$\{\abs w > \rho\}$, yang diuraikan pada ragam Fourier pada
lingkarannya, bersifat tak negatif --- jadi teorema luasnya merupakan
identitas $L^2$ beserta sebuah tanda. Sedangkan selebihnya adalah ketaksamaan itu
yang diangkut: sebab sebuah akar kuadrat (pertanyaan 3) mengubahnya menjadi
$\abs{a_2} \leq 2$; lalu sebuah pencerminan Möbius dari sebuah nilai
terlewat (pertanyaan 6) mengubah $\abs{a_2} \leq 2$ menjadi teorema
seperempatnya; lalu automorfisma cakramnya (pertanyaan 10) menyebarkan $\abs{a_2}
\leq 2$ ke seluruh cakramnya sebagai teorema kesesatannya;
lalu pengintegralan radialnya mengubah kesesatan menjadi pertumbuhan, dan
pertumbuhan menjadi penyelimutan. Dan pada setiap tahapnya kasus kesamaannya
merambat pula, yang selalu mendarat pada fungsi Koebe yang
dirotasikan --- yakni satu keluarga ekstremal bagi seluruh teorinya, persis
seperti rotasi merupakan satu-satunya ekstremal lema
Schwarz. Jadi satu ketaksamaan \omterm{def:b3:topology:interior}{interior}, ditambah kekakuan
kasus kesamaannya, mengatur seluruh geometrinya: sehingga \omterm{def:b3:conformal:conformal}{pemetaan
konformal} adalah seni memanfaatkan ketaksamaan semacam itu.

\textbf{23.} Teorema pertumbuhannya membatasi $\abs f \leq
\frac{r}{(1-r)^2}$ secara seragam pada setiap $\bar D(0, r)$ dengan $r <
1$, bagi semua $f \in \mathcal S$ sekaligus: jadi terbatas lokal, sehingga
$\mathcal S$ merupakan keluarga normal
(\cref{thm:b3:conformal:montel}). Lalu jika $f_n \in \mathcal S \to
f$ secara seragam lokal: maka $f$ \omterm{def:b3:holomorphic:holo}{holomorfik} dengan $f_n' \to f'$
secara seragam lokal (menurut Weierstrass), sehingga $f(0) = 0$ dan $f'(0) = 1$
--- jadi khususnya $f$ tak konstan --- dan Hurwitz
(\cref{exo:b3:residues:8}) membuat limit pemetaan injektifnya
injektif: sehingga $f \in \mathcal S$. Sedangkan sebuah fungsional \omterm{def:b3:topology:continuity}{kontinu} (misalnya
$f \mapsto \abs{a_2} = \frac{\abs{f''(0)}}2$) pada sebuah
kelas \omterm{def:b3:topology:compact}{kompak} \emph{mencapai} supremumnya: sehingga fungsi
ekstremalnya ada sebelum orang tahu apa itu --- yakni titik
awal setiap serangan variasi atas soal
koefisien, termasuk soal Bieberbach.

\textbf{24.} Tulis $g(w) = w + A_2w^2 + O(w^3)$ lalu komposisikan:
\[
z = g(f(z)) = f(z) + A_2f(z)^2 + O(z^3)
= z + (a_2 + A_2)z^2 + O(z^3),
\]
sehingga $A_2 = -a_2$ dan $\abs{A_2} \leq 2$ (menurut Bieberbach), dengan
kesamaannya jika dan hanya jika $\abs{a_2} = 2$, yakni jika dan hanya jika $f$ merupakan fungsi
Koebe yang dirotasikan (pertanyaan 16) --- dan lalu $g$ merupakan
balikan yang bersesuaian, yang terdefinisi pada bidang teririsnya.

\textbf{25.} Berlaku $f(z_1) - f(z_2) = (z_1 - z_2)\bigl(1 + a(z_1 +
z_2)\bigr)$. Jika $\abs a \leq \frac12$: maka untuk $z_1 \neq z_2$ di
$\mathbb D$, $\abs{a(z_1 + z_2)} < 2\abs a \leq 1$ (secara sejati, sebab
$\abs{z_1 + z_2} < 2$), sehingga faktor keduanya tak dapat lenyap:
jadi injektif, dan $f \in \mathcal S$ (sebab penormalannya sudah
tertanam). Sedangkan jika $\abs a > \frac12$: titik $s =
-\frac1a$ mempunyai $\abs s < 2$, sehingga $z_{1,2} = \frac s2 \pm
\varepsilon$ terletak di $\mathbb D$ untuk $\varepsilon > 0$ yang kecil,
berbeda, dan $z_1 + z_2 = s$ membunuh faktornya: sehingga $f(z_1)
= f(z_2)$, jadi tak injektif. Karena itu $\mathcal S$ memuat $z +
az^2$ tepat untuk $\abs a \leq \frac12$. Jadi batas Bieberbach
$\abs{a_2} \leq 2$ sangat tak terjenuhkan oleh
polinomial berderajat $2$ --- sebab koefisien fungsi Koebe
$a_n = n$ datang dari sebuah deret tak berhingga yang bersekongkol
sepanjang sinar terlewatnya, yakni perilaku yang tak dapat ditiru polinomial
mana pun (yang termasuk $\mathcal S$ hanya dengan koefisien yang mungil).

\end{solution}
